\subsection*{Lift T: controlled extraction and modeled replacement}
\label{supp:lift-causal}
For this Lift T instance, changing one extraction preference while
holding the solved game and fixed-point memory fixed produces a strict,
cooperative original-contract controller. A subsequent constructed
execution uses the checked controller after an environment-model
weakening. Both are followups designed after inspecting the baseline
failure, not held-out validation or a general synthesis correction.


\paragraph{Baseline and fixed-point agreement.}
A fresh run reproduces all 23 saved semantic artifacts of the released
baseline, including its complete 794-state, 13,874-edge native graph and
154-state, 1,170-edge original-contract view. Both graph auditors again
accept strictness and reject cooperation. A behavior-preserving observer
uses the same model and memory objects; complete graph equality,
semantic equality of the saved model and memory predicates, and
unchanged variable order establish observation parity.

A separate finite-set implementation evaluates the inspected raw
four-fixed-point recurrence on all 2,048 transformed typed states.
It uses the saved transition predicates and their original move
quantifiers, without an \emph{Alive} restriction or justice normalization.
Its region contains all 2,048 states, exactly matching the native region.
The effective 12 $X$ slices, 12 $Y$ slices, three $Z$ components and final
winning set also agree on all 57,344 state memberships.
The recurrence uses Python integer sets; a shared saved-BDD decoder
only reads predicate membership. These checks establish numerical
agreement on this input, not compiler independence, the raw core's
semantic soundness premise, or a general progress-memory theorem.

\paragraph{The transformed strict game permits the dead end.}
Node 680 belongs to the raw region and satisfies the original
response condition $C$. A focused check of its saved residual strategy
also witnesses membership in the actual transformed strict-winning
region $W_R(\cdot,0)$. All 16 legal next inputs flip the helper;
none can keep it false. The residual graph has 321 states and
5,136 edges. At every reached node it covers exactly every legal
input with a unique safe response, retaining all monitors and private
memory. Every successor has the helper true and violates the added
environment justice, so every infinite legal play satisfies the strict
liveness implication. This uses the complete saved relations, rather
than assuming that the syntactic arena matches the complete-literal
theorem. The check reuses the saved-BDD decoder; it performs no new
native generation or region solution.
Thus the game itself permits this non-cooperative behavior. The paired
intervention below shows that the same upstream data also permit an
accepted output; it does not identify extraction as the unique cause.

\paragraph{Where the continuable alternative is lost.}
Here the original environment has no justice, and
$\mathit{Alive}(s)=\exists e':\rho_e(s,e')$ equals the current one-floor
condition. The trace records two first creations of bad nodes and two
later reuses. At first creation, rule~a's raw successors have six floor
patterns, including the live $001$ and $100$ alternatives. After its
nearest-$Y$ choice, the final five-pattern candidate still contains
$100$ but not $001$; \code{satOne} selects the dead $011$.
Later rule~a prefers the existing dead node before the final choice.
Thus the source-level $001$ alternative above and the final-stage $100$
alternative have different roles. The trace contains no dropped events
or predicate references; it does not by itself certify a fair
continuation after any counterfactual choice.

\paragraph{Single-axis intervention and complete-controller checks.}
Immediately before transition \code{satOne}, the experimental constructor
intersects its final candidate with an actual helper-false input-availability
preference. It uses this intersection when nonempty, otherwise retaining
the native candidate. With helper $h$, the preference is
\[
 h(s)\ \lor\
 \exists e'\,[D_e(e')\land
       \rho_e^{\mathrm{native}}(s,e')[h=0,h'=0]] .
\]
Without a helper it uses the supplied native environment relation.
The helper-false relation is not universally identical to the original
environment relation: Section~N gives a static translation counterexample.
Lift's complete original-contract graph checks below establish its
instance-specific result independently of that universal interpretation.
This leaves helper-true candidates available. It changes no source
clause, initializer, solver, stored $X/Y/Z$ set, variable order, rank
rule, or earlier existing-state preference. There is no hardcoded floor
constraint or node ID.

One solver call supplies the same model and memory objects to the
released and experimental extractions. Their initial, transition,
ordered-justice and effective memory predicates remain semantically
equal; the baseline again matches the saved graph. Both complete outputs
are checked by the production and separately implemented graph auditors,
whose mathematical checks are unchanged and whose metadata gates
explicitly identify the nonrelease variant.
The production reader reconstructs saved predicates as BDDs; the other
auditor evaluates the exported decision DAG with a separate reader and
compares it to separately encoded source equations. Each run checks all
$1{,}048{,}576$ original current/next state pairs for correspondence and
$2{,}464$ typed input obligations on its 154-node original view.
The second auditor imports neither the production checker nor its
certificate and uses a separate reachability-based fairness test.
Both still depend on the same compiler exports, typing information,
and recorded translation convention; implementations were mutually
visible. The shared decoder limitations for the fixed-point and
residual-strategy checks above are distinct from this graph-audit pair.
Table~\ref{tab:lift-choice-intervention} reports the original-contract
views. The experimental full native graph has 778 states and 13,624
edges; 726 of its 13,624 final candidates are refined, with no
empty-preference fallback. This is one code-level change affecting many
choices, not four patched witness edges.

\begin{table}[htbp]
\centering\small
\caption{Paired extraction results on Lift T and six prespecified
successful controls. Each cell gives original-view states/edges and
strict/cooperative verdicts; P/F mean pass/fail. Refinements count
candidate sets changed by the preference. These are overlapping
followups, not additional members of the eleven-candidate cohort.}
\label{tab:lift-choice-intervention}
\begin{tabular}{@{}lccc@{}}
\toprule
Input & Released & Experimental & Refinements\\
\midrule
Lift T & 154/1170, P/F & 154/1176, P/P & 726\\
Arbiter T & 56/448, P/P & 56/448, P/P & 0\\
Minepump T1 & 12/40, P/P & 12/40, P/P & 0\\
Minepump T2 & 4/16, P/P & 4/16, P/P & 0\\
Traffic S T & 12/80, P/P & 12/80, P/P & 0\\
Traffic U P & 68/323, P/P & 68/323, P/P & 0\\
Traffic U T & 56/416, P/P & 56/416, P/P & 0\\
\bottomrule
\end{tabular}
\end{table}

The six controls were declared before their new executions and retain
both verdicts with zero refinements and zero fallbacks. They establish
preservation of those successes when the preference is dormant, not
generality when it is active. For Lift, unchanged game, raw region and
memory suffice for this constructor to return a checked cooperative
view after changing concrete extraction choices. One-step availability alone is no
fairness proof; the positive conclusion comes from complete-controller
checks. No sole-defect, missing-$W$ cause, or general synthesis-repair
claim follows.

\paragraph{A constructed environment change and exact handoff.}
The old contract is Lift T with the additional environment restriction
that all current and next buttons are false. A source-coded old
controller stays at floor $100$ and performs the actual response-monitor
updates. Its two reachable cancellation states and four edges pass
strictness and cooperation under this maintenance contract.
Two modeled plants execute four old-controller steps, with next
cancellation values $0,1,1,0$. Removing only the maintenance restriction
then exposes the original Lift T environment. Both plants pause at
the existing state $b=000,c=0,f=100,m=000$.

After this change, a fresh paired native run generates both candidates,
followed by four fresh graph judgments; no cached candidate is installed.
The fixed order is released controller, then experimental controller.
One admission policy accepts the first strict candidate; the other
requires both strictness and cooperation. Each also requires exact
agreement between the current full source valuation and an allowed
initial native node. Installation preserves every source field and
logical monitor and initializes the new controller's private memory to
that matching node. The handoff therefore uses an already
initial-compatible paused state; it is not arbitrary-state migration
or a reset of pending requests.

\paragraph{Executed outcomes.}
Both arms execute the five inputs of the known failing prefix,
in $(b,c)$ notation:
$(010,1),(001,1),(100,0),(101,0),(101,0)$.
The strict-only policy installs the released controller and reaches
floor $011$ after these five post-change steps. Independent enumeration
finds all 16 possible original next inputs illegal.
The strict-and-cooperative policy rejects that candidate, installs the
experimental one, and reaches floor $100$ instead.
It then executes the lexicographically smallest legal input at each
step, with false ordered before true in $(b_1,b_2,b_3,c)$.
Within the declared 32-step continuation bound, four more transitions
give native nodes
$680\to524\to544\to434\to434$.
The final self-loop has $b=000,c=0,f=001,m=000$, meeting all three
system justices; there are no original environment justices.
Native IDs are local to each graph: experimental node 680 is not the
released graph's dead node 680.

A separate no-replacement counterfactual applies $(010,0)$ twice to the
stationary old controller under the weakened contract. The first step
is legal; the second proposed response violates the movement guarantee
and is rejected. Thus keeping this old controller does not solve the
changed model.

The independent episode audit imports neither the episode runner nor
its transition implementation. It checks the 47-event log, source
semantics, exact handoff, installed-controller hashes and actual native
responses, and binds the four fresh full-graph judgments. It performs
24 environment and 24 system source checks, including the rejected
counterfactual step, and 42 native-edge/response checks. All six
supervised processes, comprising the outer execution, one native pair
and four judges, complete with verified cleanup.
The executed loop witnesses one infinite continuation; the separate
complete-graph certificates justify the controller properties.
This is a constructed in-process episode with a known-failure prefix,
not hardware deployment, newly learned repair, arbitrary memory
transfer, or a runtime-cost benchmark.

\paragraph{Bounds, attempts, and provenance.}
Table~\ref{tab:lift-followup-bounds} gives the declared successful-run
limits. Stream bounds apply separately to stdout and stderr.
RSS caps apply to the owned process groups, not a measured aggregate
tree; output bounds are sampled termination thresholds with final
checks, not exact quotas. The episode's outer final \code{wait4}
resource value can include waited descendants on this host.
Native and checker children retain a 16-MiB per-artifact limit; the
outer episode's per-file limit is 1~GiB so it can capture the already
bounded native streams. Each launch requires 1~GiB of free disk space.

\begin{table}[htbp]
\centering\small
\caption{Bounds for the accepted diagnostic and episode runs.
The episode includes serial child processes with their own limits.}
\label{tab:lift-followup-bounds}
\begin{tabular}{@{}lrrrr@{}}
\toprule
Process & Timeout (s) & RSS (GiB) & Stream (MiB) & Output (MiB)\\
\midrule
Fresh baseline & 10 & 1 & 256 & 512\\
Observer / paired extraction & 30 & 1 & 512 & 768\\
Finite-set / memory comparison & 120 & 1 & 1 & 32\\
Production graph judge & 60 & 1 & 8 & 64\\
Independent graph judge & 180 & 1 & 8 & 64\\
Outer modeled episode & 600 & 2 & 8 & 1024\\
\bottomrule
\end{tabular}
\end{table}

\paragraph{Observed durations of the completed episode.}
Table~\ref{tab:lift-episode-observed-times} extracts recorded
\code{elapsed\_seconds} from the successful episode's process receipts.
These are supervised wall durations including capture, finalization,
and I/O; the timeout column is the declared cap, not an observation.
The outer run contains the paired native generation and four serial
graph judgments, together with modeled steps and orchestration. Its
28.36~s must not be added to those child durations. The separate saved-event
replay occurred later and is outside that outer run. No solver-only or
extraction-only time is inferred from the paired process.

\begin{table}[htbp]
\centering\small
\caption{Recorded process durations for one completed constructed
Lift T episode. Nested rows are components of the outer run;
these observations are not a latency or scaling benchmark.}
\label{tab:lift-episode-observed-times}
\begin{tabular}{@{}lrr@{}}
\toprule
Process & Observed wall (s) & Timeout (s)\\
\midrule
Outer modeled episode & 28.36 & 600\\
\quad Paired native generation & 10.82 & 30\\
\quad Released: production judge & 0.12 & 60\\
\quad Released: separate judge & 8.60 & 180\\
\quad Experimental: production judge & 0.15 & 60\\
\quad Experimental: separate judge & 8.35 & 180\\
Later saved-event replay & 0.11 & 120\\
\bottomrule
\end{tabular}
\end{table}

These measurements describe one completed run of the three-floor
example after the three retained technical UNKNOWN attempts below.
They are not a statistical sample, online adaptation latency,
hardware-normalized comparison, worst-case guarantee, or a scaling
result beyond the bounded models. Process-group RSS caps remain limits;
they are not measured aggregate-tree memory consumption.

The first observer attempt exceeded its 256-MiB stdout limit and
remains UNKNOWN; a separately recorded followup raised only stream
and output-tree limits to 512 and 768~MiB. The episode likewise retains
three earlier attempts as technical UNKNOWNs. Attempt~01 stopped when
the outer single-child supervisor detected a nested driver; the
followup called the existing driver mains in the episode process while
each native/checker child retained its own supervisor.
Attempt~02's stream capture hit the inherited outer 16-MiB file limit;
the next attempt raised that outer limit to 1~GiB.
Attempt~03 produced the complete event log and five valid inner
processes, but its final \code{wait4} RSS value of 557,793,280 bytes
exceeded the outer 512-MiB cap, so it remains resource UNKNOWN.
Attempt~04 raised the outer RSS cap to 2~GiB, retained all inner and
output limits, and completed; its independent episode audit passes.
Scientific inputs, candidate order, handoff, common input prefix and
acceptance criteria stayed fixed. Each followup has a separate plan,
output directory and receipt; no automatic retry rewrites an attempt.

The author-side records bind source, class, model, graph, checker,
process and event-log identities. Their entry points are
\path{CAUSAL_COMPARISON_CLOSURE.json},
\path{SUCCESS_CONTROLS_RESULT.json}, and
\path{adaptation_episode_04/EXECUTION_RECEIPT.json}, followed by
\path{adaptation_episode_04/independent/AUDIT_RESULT_attempt01.json}.
These distinguish the original released result, experimental result
and executed episode. The fourth public archive in Section~K
preserves them with explicit original/copy identities. Its consumer
replays saved semantics and the event audit; it does not regenerate
the native producer or execute a new adaptation episode.
